\chapter{Craniotomy}

\section*{What Is This Surgery?}
A craniotomy is a major neurosurgical procedure involving the temporary removal of a section of the skull, known as a bone flap, to gain direct access to the brain and its surrounding structures. This intricate operation is performed to diagnose and treat a wide array of intracranial pathologies, ranging from life-threatening traumatic injuries and vascular malformations to benign and malignant brain tumors. The primary objective is to safely expose the affected area of the brain, allowing for precise surgical intervention under direct visualization, often with the aid of a surgical microscope or endoscope. Upon completion of the intracranial work, the bone flap is meticulously resecured to restore the integrity and protective function of the cranial vault. This procedure is critical for alleviating intracranial pressure, excising pathological tissue, repairing vascular anomalies, or draining fluid collections that threaten neurological function and survival.

\section*{Alternative Names}
\begin{itemize}
  \item Open Brain Surgery
  \item Intracranial Surgery
  \item Skull Opening Procedure
  \item Bone Flap Craniotomy
  \item Cranial Trephination (historically, though modern craniotomy is more precise and involves bone flap replacement)
  \item Specific approaches: Pterional Craniotomy, Suboccipital Craniotomy, Frontotemporal Craniotomy, Orbitozygomatic Craniotomy
\end{itemize}

\section*{Relevant Surgical Anatomy}
A thorough understanding of cranial and intracranial anatomy is paramount for safe and effective craniotomy. The skull, composed of several fused bones (frontal, parietal, temporal, occipital, sphenoid, ethmoid), provides a rigid protective casing for the brain. Key surgical landmarks on the skull surface include the pterion (a weak point where frontal, parietal, temporal, and sphenoid bones meet, often targeted for middle fossa access), asterion (junction of parietal, temporal, and occipital bones), and superior temporal line. Beneath the skull lie the meninges: the dura mater (a tough, fibrous outer layer), the arachnoid mater (a delicate, web-like middle layer), and the pia mater (a thin layer intimately adhering to the brain surface). The subarachnoid space, between the arachnoid and pia, contains cerebrospinal fluid (CSF) and major cerebral blood vessels.

The brain itself is divided into cerebrum, cerebellum, and brainstem, each with distinct functions and vascular territories. The cerebrum comprises frontal, parietal, temporal, and occipital lobes, with critical functional areas such as the motor cortex, sensory cortex, language centers (Broca's and Wernicke's areas), and visual cortex. Major arterial supply to the brain is derived from the internal carotid arteries and vertebral arteries, forming the Circle of Willis at the base of the brain. Key vessels encountered during craniotomy include the Middle Cerebral Artery (MCA) and its branches (often involved in aneurysms), Anterior Cerebral Artery (ACA), and Posterior Cerebral Artery (PCA). Venous drainage occurs via superficial cortical veins draining into the superior sagittal sinus, transverse sinuses, and sigmoid sinuses, ultimately leading to the internal jugular veins. Cranial nerves emerge from the brainstem and base of the skull, and their preservation is critical; specific nerves (e.g., optic, oculomotor, trigeminal, facial, vestibulocochlear) are at risk depending on the surgical approach and pathology location. Lymphatic drainage of the brain is still being elucidated but primarily involves perivascular spaces and drainage into cervical lymph nodes.

\section{Surgical Principle \& Rationale}
The fundamental principle of a craniotomy is to create a controlled, temporary opening in the skull to allow direct surgical access to the intracranial contents while minimizing trauma to surrounding healthy brain tissue and preserving cranial integrity. This is achieved by carefully elevating a section of the skull bone, known as the bone flap, which is then replaced at the end of the procedure. The rationale is to provide a wide enough surgical corridor for the neurosurgeon to safely manipulate delicate brain structures, resect lesions, clip aneurysms, or evacuate hematomas, often under high magnification using a surgical microscope or endoscope. This direct visualization and access are crucial for precise intervention in a confined and vital anatomical space.

The technical goals of a craniotomy are multifaceted. Firstly, to achieve adequate exposure of the pathology while protecting eloquent brain regions and critical neurovascular structures. Secondly, to ensure meticulous hemostasis throughout the procedure to prevent intraoperative and postoperative hemorrhage. Thirdly, to perform the specific surgical task (e.g., tumor resection, aneurysm clipping, hematoma evacuation) with maximal efficacy and minimal collateral damage. Finally, to restore the protective barrier of the skull and achieve a watertight dural closure to prevent cerebrospinal fluid leakage and infection, thereby facilitating optimal neurological recovery.

\section{History \& Surgical Evolution}
The concept of opening the skull has ancient roots, with evidence of trepanation dating back thousands of years across various cultures, often performed for ritualistic purposes, to relieve headaches, or to treat skull fractures. Early trepanations were crude, often resulting in high mortality, but some patients survived, indicating a rudimentary understanding of the brain's role in health. Hippocrates, in ancient Greece, described trephination for head injuries.

Modern craniotomy began to take shape in the late 19th century with the advent of antiseptic surgical techniques pioneered by Joseph Lister. This allowed for safer intracranial procedures by drastically reducing infection rates. Key figures in the development of neurosurgery include William Macewen, a Scottish surgeon who, in the 1880s, successfully drained brain abscesses and removed brain tumors, demonstrating the feasibility of intracranial surgery. Harvey Cushing, an American surgeon, is widely regarded as the father of modern neurosurgery. Beginning in the early 20th century, Cushing revolutionized the field by emphasizing meticulous hemostasis, developing specialized instruments, and establishing the first neurosurgical residency program. His work significantly reduced the mortality and morbidity associated with brain surgery, transforming it from a desperate measure into a viable therapeutic option. Subsequent advancements in imaging (X-rays, angiography, CT, MRI), microsurgical techniques, neuroanesthesia, intraoperative neurophysiological monitoring, and neuronavigation have further refined craniotomy, making it safer, more precise, and applicable to an ever-wider range of complex intracranial pathologies.

\section{Clinical Indications}
A craniotomy is indicated for a variety of critical neurological conditions that require direct surgical intervention to diagnose, treat, or alleviate symptoms. Specific indications include:

\begin{itemize}
  \item Resection of primary or metastatic brain tumors, including gliomas, meningiomas, pituitary adenomas, and acoustic neuromas, to achieve gross total resection or debulking.
  \item Clipping of cerebral aneurysms, both ruptured (to prevent re-bleeding and manage subarachnoid hemorrhage) and unruptured (to prevent initial rupture, especially for larger or symptomatic aneurysms).
  \item Excision of arteriovenous malformations (AVMs) and cavernous malformations to prevent hemorrhage, control medically refractory seizures, or alleviate mass effect.
  \item Evacuation of traumatic intracranial hematomas, such as epidural, subdural, or intraparenchymal hematomas, to relieve mass effect, brain compression, and impending herniation.
  \item Drainage and excision of brain abscesses or other intracranial infections that are refractory to antibiotic therapy or cause significant mass effect.
  \item Decompression of cranial nerves, such as in trigeminal neuralgia or hemifacial spasm, through microvascular decompression to relieve neurovascular compression.
  \item Resection of epileptogenic foci in patients with medically refractory epilepsy, such as temporal lobectomy or lesionectomy.
  \item Repair of cerebrospinal fluid (CSF) leaks that are resistant to conservative management, often following trauma or previous surgery.
  \item Biopsy of deep-seated brain lesions that cannot be diagnosed through less invasive means, providing tissue for definitive pathological diagnosis.
  \item Evacuation of spontaneous intracerebral hemorrhage, particularly large or superficial clots causing significant neurological deficit or mass effect.
\end{itemize}

\section*{Clinical Positioning}
Craniotomy often serves as a primary, definitive operative procedure for many acute and chronic intracranial diseases. For life-threatening conditions such as ruptured cerebral aneurysms, large traumatic hematomas causing significant mass effect, or rapidly growing malignant brain tumors, craniotomy is frequently the first-line treatment to save life or preserve neurological function. In these acute scenarios, the immediate and direct access to the pathology offered by craniotomy is paramount, as delays can lead to irreversible brain damage or death.

However, craniotomy can also be a secondary or salvage procedure. This may occur when less invasive treatments, such as endovascular coiling for aneurysms, stereotactic radiosurgery for small tumors, or medical management for infections, have failed, are contraindicated, or are not feasible. For instance, a craniotomy might be performed to clip an aneurysm that could not be successfully coiled due to anatomical constraints, or to resect a tumor that recurred after radiation therapy. It can also be an adjunct to other therapies, such as when a biopsy via craniotomy is needed to guide subsequent chemotherapy or radiation, or when it is part of a multi-stage treatment plan for complex vascular malformations. The decision to perform a craniotomy, and its role as primary or secondary, is highly dependent on the specific pathology, its location, the patient's clinical status, the availability of alternative treatments, and the overall risk-benefit profile.

\section*{Surgical Technique \& Procedure}
Craniotomy is performed under general anesthesia, ensuring the patient is completely unconscious and pain-free throughout the extensive procedure. Patient positioning is meticulously planned and executed, often utilizing a skull clamp (e.g., Mayfield clamp) to rigidly fix the head in a specific orientation that optimizes surgical access while protecting the airway and preventing pressure injuries. The exact position (e.g., supine, lateral, prone, park bench) depends on the location of the lesion and the chosen surgical approach. The scalp is prepared by clipping or shaving hair, followed by thorough antiseptic cleansing with agents like povidone-iodine or chlorhexidine.

The surgical steps typically involve:
\begin{itemize}
  \item \textbf{Scalp Incision and Flap Creation:} A curvilinear, horseshoe, or question-mark shaped incision is made in the scalp, carefully designed to maximize exposure while preserving blood supply to the scalp flap. The scalp and underlying muscle layers (e.g., temporalis muscle for pterional approach) are then carefully dissected and reflected to expose the skull.
  \item \textbf{Burr Hole Placement:} Multiple small holes, typically 3-5, are drilled into the skull using a high-speed drill with a safety stop to prevent dural injury. The number and precise placement of these burr holes define the perimeter of the bone flap, often guided by neuronavigation systems.
  \item \textbf{Bone Flap Creation:} A specialized surgical saw called a craniotome is used to connect the burr holes, creating a precise cut through the full thickness of the skull bone. The bone flap is then carefully elevated, often hinged on a muscle pedicle to preserve blood supply, or completely removed and placed in a sterile basin for later replacement.
  \item \textbf{Dural Opening:} The dura mater, the tough outer membrane covering the brain, is meticulously incised and reflected, typically in a curvilinear or C-shaped fashion. This step requires extreme care to avoid injury to the underlying brain or bridging veins, which can cause significant hemorrhage.
  \item \textbf{Intracranial Work:} Under microscopic or endoscopic guidance, the neurosurgeon performs the specific procedure, such as tumor resection (using ultrasonic aspirators, bipolar cautery), aneurysm clipping (using titanium clips), hematoma evacuation, or microvascular decompression. Intraoperative neurophysiological monitoring (e.g., motor evoked potentials, somatosensory evoked potentials) may be used to protect eloquent brain regions.
  \item \textbf{Dural Closure:} Once the intracranial work is complete, meticulous hemostasis is achieved. The dura mater is then carefully closed, often in a watertight fashion using fine non-absorbable sutures, to prevent cerebrospinal fluid leakage. A dural substitute (e.g., synthetic graft, autologous fascia) may be used if primary closure is not possible or to reinforce the repair.
  \item \textbf{Bone Flap Replacement and Fixation:} The bone flap is precisely repositioned and secured to the surrounding skull using small titanium plates and screws, ensuring stability, proper alignment, and a good cosmetic outcome.
  \item \textbf{Scalp Closure:} The muscle and scalp layers are meticulously closed in layers with absorbable and non-absorbable sutures. A subgaleal drain may be placed beneath the scalp to prevent fluid accumulation (hematoma or seroma). A sterile dressing is applied.
\end{itemize}

\section*{Preoperative Workup \& Preparation}
Thorough preoperative preparation is crucial for optimizing patient safety and surgical outcomes in craniotomy. This typically involves a comprehensive medical and neurological evaluation to assess the patient's overall health and the specific intracranial pathology.

Key preparatory steps include:
\begin{itemize}
  \item \textbf{Detailed Neurological Examination:} To establish a baseline of neurological function, identify any existing deficits, and monitor for changes.
  \item \textbf{Laboratory Tests:} Complete blood count (CBC), electrolyte panel, renal and hepatic function tests, coagulation studies (PT/INR, PTT), and blood type and cross-match are essential.
  \item \textbf{Advanced Imaging:} High-resolution magnetic resonance imaging (MRI) with and without contrast, computed tomography (CT) scans, CT angiography (CTA), magnetic resonance angiography (MRA), or digital subtraction angiography (DSA) are often performed to precisely localize the lesion, characterize its nature, and plan the surgical approach. Functional MRI (fMRI) or diffusion tensor imaging (DTI) may be used to map critical brain areas (e.g., motor cortex, language centers) to guide surgical planning and minimize deficits.
  \item \textbf{Medical Clearance:} Comprehensive cardiac and pulmonary evaluations, often including an electrocardiogram (ECG), chest X-ray, and sometimes echocardiography or pulmonary function tests, are obtained to ensure the patient is medically fit for general anesthesia and major surgery. Anesthesia consultation is mandatory to assess anesthetic risks and plan perioperative management.
  \item \textbf{Medication Management:} Anticoagulant and antiplatelet medications (e.g., warfarin, aspirin, clopidogrel) must be discontinued several days to a week prior to surgery, often with bridging therapy if necessary, to minimize bleeding risk. Steroids (e.g., dexamethasone) may be initiated to reduce cerebral edema, and anticonvulsants may be prescribed prophylactically to prevent postoperative seizures.
  \item \textbf{NPO Status:} Patients are instructed to be NPO (nil per os - nothing by mouth) for a specified period (typically 6-8 hours) before surgery to prevent aspiration during anesthesia induction.
  \item \textbf{Informed Consent:} A detailed discussion with the patient and family regarding the procedure, its risks, benefits, alternatives, and expected recovery is conducted, and informed consent is obtained.
  \item \textbf{Preoperative Antibiotics:} Intravenous antibiotics (e.g., cefazolin) are administered shortly before the skin incision to minimize the risk of surgical site infection.
\end{itemize}

\section*{Contraindications \& Precautions}
\textbf{Absolute Contraindications:}
\begin{itemize}
  \item \textbf{Severe, Uncontrolled Coagulopathy:} An inability to correct significant bleeding disorders (e.g., severe hemophilia, severe thrombocytopenia) poses an unacceptable risk of catastrophic intracranial hemorrhage.
  \item \textbf{Active Systemic Infection or Sepsis:} Performing elective craniotomy in the presence of widespread infection significantly increases the risk of surgical site infection, meningitis, or brain abscess.
  \item \textbf{Irreversible Brain Damage or Futility:} In cases of severe, irreversible brain injury (e.g., brain death, devastating anoxic injury) where there is no reasonable expectation of meaningful neurological recovery, craniotomy may be deemed futile.
  \item \textbf{Severe, Uncorrectable Medical Comorbidities:} Patients with end-stage organ failure (e.g., severe heart failure, respiratory failure, liver failure) that cannot be optimized may not tolerate the physiological stress of surgery and anesthesia.
\end{itemize}

\textbf{Relative Contraindications and Precautions:}
\begin{itemize}
  \item \textbf{Extreme Advanced Age with Multiple Comorbidities:} While age alone is not a contraindication, very elderly patients with significant medical burdens (e.g., severe cardiovascular disease, frailty) may have increased surgical and anesthetic risks, requiring careful risk-benefit assessment.
  \item \textbf{Prior Radiation to the Surgical Field:} Previous radiation therapy can alter tissue planes, increase vascularity, impair wound healing, and make subsequent surgery more challenging and risky due to tissue fibrosis and necrosis.
  \item \textbf{Severe Cerebral Edema with Impending Herniation:} In some acute cases, a decompressive craniectomy (where the bone flap is not immediately replaced) may be a safer initial approach than a craniotomy with immediate bone flap replacement, to allow for brain swelling.
  \item \textbf{Poor Neurological Status with Borderline Prognosis:} The potential benefits of surgery must be carefully weighed against the risks in patients with very poor preoperative neurological function (e.g., GCS < 8 without reversible cause), as the likelihood of meaningful recovery may be low.
  \item \textbf{Patient Refusal:} A competent patient's refusal of surgery, after thorough informed consent, is a contraindication.
  \item \textbf{Uncontrolled Hypertension or Diabetes:} These conditions must be meticulously optimized preoperatively to minimize intraoperative and postoperative complications such as hemorrhage, stroke, or impaired wound healing.
  \item \textbf{Active Substance Abuse:} May complicate anesthesia, pain management, and postoperative recovery.
\end{itemize}

\section*{Complications \& Risks}
Craniotomy, as a major neurosurgical procedure, carries inherent risks, some of which are general to any surgery and others specific to brain intervention.

\begin{itemize}
  \item \textbf{General Surgical Risks:}
    \begin{itemize}
      \item \textbf{Bleeding:} Intraoperative or postoperative hemorrhage, which can lead to hematoma formation (epidural, subdural, intracerebral) requiring urgent reoperation and potentially causing neurological deterioration.
      \item \textbf{Infection:} Surgical site infection (wound infection, osteomyelitis of the bone flap), meningitis (infection of the brain coverings), or brain abscess, which can be life-threatening.
      \item \textbf{Anesthesia Complications:} Allergic reactions to medications, cardiac events (myocardial infarction, arrhythmia), respiratory depression, pneumonia, deep vein thrombosis (DVT), pulmonary embolism, and malignant hyperthermia.
    \end{itemize}
  \item \textbf{Procedure-Specific Neurological Risks:}
    \begin{itemize}
      \item \textbf{Neurologic Deficit:} New or worsened weakness, paralysis (hemiparesis/hemiplegia), speech difficulties (aphasia), vision changes (hemianopsia), cognitive impairment, or sensory loss, depending on the brain area involved and the extent of manipulation.
      \item \textbf{Seizures:} Can occur immediately postoperatively or develop later, sometimes requiring long-term anticonvulsant medication.
      \item \textbf{Cerebral Edema/Brain Swelling:} Swelling of the brain tissue, which can increase intracranial pressure and potentially lead to herniation, requiring medical management or further surgical decompression.
      \item \textbf{Cerebrospinal Fluid (CSF) Leak:} Leakage of CSF from the wound, nose (rhinorrhea), or ear (otorrhea), increasing the risk of meningitis and often requiring surgical repair.
      \item \textbf{Stroke:} Ischemic stroke (due to vessel occlusion or vasospasm) or hemorrhagic stroke (due to bleeding) during or after surgery, leading to permanent neurological deficits.
      \item \textbf{Hydrocephalus:} Accumulation of CSF in the brain ventricles, potentially requiring a shunt placement (ventriculoperitoneal shunt) to drain excess fluid.
      \item \textbf{Cranial Nerve Injury:} Damage to specific cranial nerves, leading to deficits such as facial weakness (CN VII), hearing loss (CN VIII), double vision (CN III, IV, VI), or swallowing difficulties (CN IX, X).
      \item \textbf{Venous Air Embolism:} Air entering the venous system, particularly in procedures performed with the patient in a sitting position, which can be life-threatening.
      \item \textbf{Bone Flap Complications:} Infection of the bone flap, non-union, resorption, malposition, or cosmetic deformity, sometimes requiring revision surgery.
      \item \textbf{Vascular Injury:} Damage to major intracranial arteries or veins, potentially leading to hemorrhage, stroke, or pseudoaneurysm formation.
    \end{itemize}
\end{itemize}

\section*{Postoperative Recovery Timeline}
Immediately following a craniotomy, the patient is transferred to the Post-Anesthesia Care Unit (PACU) for initial recovery from anesthesia. During this critical period, vital signs are continuously monitored, and frequent neurological checks (assessing consciousness, pupil reactivity, motor strength, and speech) are performed to detect any immediate complications such as hemorrhage or worsening cerebral edema. Pain management is initiated with intravenous analgesics, and antiemetics are given to prevent nausea and vomiting, which can increase intracranial pressure. Once stable, the patient is typically transferred to a specialized neurosurgical intensive care unit (NICU) or a dedicated neurosurgical ward for continued close observation.

In the NICU, monitoring continues with close attention to intracranial pressure (ICP) if a monitor has been placed, blood pressure control to maintain adequate cerebral perfusion, and fluid balance. Head elevation (typically 30 degrees) is usually maintained to promote venous drainage and reduce cerebral edema. Prophylactic anticonvulsants and steroids may be continued and gradually tapered. Early mobilization, often starting with sitting up and ambulation within 24-48 hours, is encouraged to prevent complications like deep vein thrombosis and pneumonia. Physical, occupational, and speech therapy evaluations may begin early to assess and address any new neurological deficits. The hospital stay typically ranges from 5 to 10 days, depending on the complexity of the surgery, the patient's recovery trajectory, and the development of any complications. Long-term recovery can involve extensive rehabilitation and may take several weeks to months for full functional improvement, with gradual resumption of normal activities.

\section*{Surgical Findings \& Pathology}
During a craniotomy, the surgeon meticulously documents the intraoperative findings, which include the appearance and consistency of the brain tissue, the nature and extent of the lesion (e.g., tumor size, location, invasiveness; aneurysm morphology; hematoma volume), its relationship to critical neurovascular structures, and the completeness of resection or repair. Digital images or video recordings are often taken to document these findings. Any tissue removed is immediately sent for pathological examination.

The pathology report, derived from the resected tissue, is often the most critical "result" for conditions like brain tumors. This report provides a definitive diagnosis, including the histological type, grade (e.g., WHO Grade I-IV for gliomas), and molecular characteristics (e.g., IDH mutation, 1p/19q co-deletion for gliomas) of the tumor, and assesses the margins of resection (e.g., positive or negative for tumor cells). These findings are paramount in guiding further treatment decisions, such as the need for adjuvant radiation therapy or chemotherapy, and for determining prognosis. For vascular lesions like aneurysms, successful clipping confirmed by intraoperative angiography or postoperative imaging indicates a reduced risk of re-rupture, and the pathology report would confirm the presence of an aneurysm sac if a sample was taken.

\section*{Expected Postoperative Course}
A normal, successful postoperative course following a craniotomy signifies resolution or significant improvement of the underlying intracranial pathology with a stable or improved neurological status and minimal complications. Patients can expect initial pain and discomfort around the incision site, which is managed with analgesics and typically improves over the first few days to weeks. Headaches are common but should gradually subside. Neurological function should either be preserved at baseline or show progressive improvement in any preoperative deficits. Swelling around the incision and potentially some facial swelling are normal and resolve over days to weeks. The surgical wound should heal cleanly, without signs of infection, excessive drainage, or CSF leakage. Patients are typically able to gradually increase their activity levels, progressing from bed rest to ambulation within days, and resuming light activities within a few weeks. The ultimate goal is to return to a pre-illness functional status or achieve the best possible neurological outcome given the underlying pathology.

\section{Postoperative Care \& Follow-up}
Post-craniotomy follow-up is essential for monitoring recovery, detecting complications, and managing long-term care.

\begin{itemize}
  \item \textbf{Neurosurgical Clinic Visits:} Regular appointments with the neurosurgeon are scheduled, typically at 2 weeks, 1 month, 3 months, and then annually, to monitor neurological recovery, assess wound healing, and address any concerns.
  \item \textbf{Suture/Staple Removal:} Scalp sutures or staples are typically removed by the neurosurgeon or a trained nurse 7-14 days post-surgery.
  \item \textbf{Serial Neuroimaging:} Follow-up MRI or CT scans are routinely performed at scheduled intervals (e.g., 3, 6, 12 months, then annually) to monitor for tumor recurrence, aneurysm stability, or other long-term changes related to the pathology or surgical site.
  \item \textbf{Neurological Rehabilitation:} Physical, occupational, and speech therapy may be continued for weeks to months, or even longer, to help patients regain strength, coordination, cognitive function, and communication skills, tailored to individual deficits.
  \item \textbf{Medication Management:} Anticonvulsant medications may be continued for a period (e.g., 3-6 months post-tumor resection) or long-term, depending on the risk of seizures. Steroids are typically tapered over several weeks.
  \textbf{Activity Resumption:} Gradual resumption of normal activities is advised, with specific restrictions on heavy lifting, strenuous exercise, and contact sports for several weeks to months to protect the healing skull. Driving restrictions are common initially, especially if seizures have occurred or if there are significant neurological deficits.
  \item \textbf{Oncology Follow-up:} For patients with malignant brain tumors, follow-up with a neuro-oncologist is crucial for ongoing management, including chemotherapy or radiation therapy, and surveillance for recurrence.
\end{itemize}
Patients are educated on warning signs such as worsening headache, fever, new neurological deficits, seizures, wound changes (redness, swelling, drainage), or excessive fatigue, and instructed to seek immediate medical attention if these occur.

\section*{Epidemiology}
The epidemiology of craniotomy is directly linked to the incidence of the underlying neurological conditions it treats. Brain tumors, both primary and metastatic, represent a significant indication, with approximately 86,000 new cases of primary brain and central nervous system tumors diagnosed annually in the United States, and metastatic brain tumors occurring in 20-40\% of cancer patients. Cerebral aneurysms are found in about 3-5\% of the general population, with roughly 30,000 ruptures occurring each year in the U.S., leading to subarachnoid hemorrhage. Traumatic brain injury (TBI) is a leading cause of craniotomy, with millions of emergency department visits annually, and a subset requiring surgical intervention for hematoma evacuation. Arteriovenous malformations (AVMs) are less common, with an estimated prevalence of 0.1\% of the population. The age and sex distribution for craniotomy vary widely depending on the specific pathology; for instance, gliomas tend to affect older adults, while AVMs are often diagnosed in younger individuals, and trauma affects all age groups. Mortality and morbidity rates are highly variable, ranging from low (e.g., <1-2\%) for elective procedures on benign, superficial lesions to significantly higher (e.g., 10-30\% or more) for emergency procedures involving large, deep-seated malignant tumors or severe traumatic brain injuries with significant mass effect. Overall, craniotomy remains a frequently performed and critical neurosurgical procedure globally, with thousands performed annually.

\section*{Outcomes \& Prognosis}
The outcomes and prognosis following a craniotomy are highly dependent on the underlying pathology, the patient's preoperative neurological status, the location and extent of the lesion, and the occurrence of postoperative complications. For acute, life-threatening conditions like traumatic epidural hematomas, craniotomy can be life-saving, with good functional recovery possible if performed promptly and effectively. For ruptured cerebral aneurysms, successful clipping prevents re-rupture, significantly improving survival, though neurological deficits from the initial hemorrhage may persist. In brain tumor surgery, the extent of resection is a critical prognostic factor; gross total resection of benign tumors often leads to cure, while for malignant tumors, it can extend survival and improve quality of life, often in conjunction with adjuvant therapies. For example, studies like the EORTC 22981/26981 trial have shown the benefit of adjuvant chemoradiation after resection of glioblastoma.

Functional outcomes vary widely. Patients may experience full recovery, while others may have persistent neurological deficits requiring long-term rehabilitation. Factors such as age, preoperative Karnofsky Performance Status (KPS), the presence of significant comorbidities, and the specific brain region involved also significantly influence prognosis. Recurrence rates are a major concern for malignant brain tumors, necessitating ongoing surveillance. Overall, advancements in surgical techniques, neuroimaging, intraoperative monitoring, and perioperative care have steadily improved the safety and efficacy of craniotomy, leading to better survival rates and functional outcomes for a broad spectrum of complex neurological conditions.

\section*{References}
\begin{enumerate}
  \item MedlinePlus. Craniotomy. U.S. National Library of Medicine; 2024. Available from: \url{https://medlineplus.gov/ency/article/002986.htm}
  \item Greenberg MS. Handbook of Neurosurgery. 9th ed. Thieme; 2020.
  \item Youmans and Winn Neurological Surgery. 8th ed. Elsevier; 2023.
  \item American Association of Neurological Surgeons (AANS). Brain Tumors. Available from: \url{https://www.aans.org/Patients/Neurosurgical-Conditions-and-Treatments/Brain-Tumors}
\end{enumerate}

\clearpage